% SPDX-FileCopyrightText: 2026 Will Cook
% SPDX-License-Identifier: CC-BY-4.0
% System-architecture paper; build with tectonic.
\documentclass[10pt]{article}
\usepackage{paper-house-style}
\usepackage{booktabs,array,float,placeins,multicol}
\usepackage[colorlinks=true,linkcolor=linkinternal,urlcolor=linkexternal,citecolor=linkinternal]{hyperref}
\hypersetup{
  bookmarksnumbered=true,bookmarksopen=true,bookmarksopenlevel=1,
  pdftitle={Problem-Sized Lean Worlds},
  pdfauthor={Will Cook},
  pdfsubject={Persistent, checkable research records for AI-assisted mathematics},
  pdflang={en-GB},
  pdfkeywords={formal mathematics, Lean 4, AI agents, research infrastructure, open problems, Erdos problems, reproducibility},
}
\newcommand{\repobase}{https://github.com/wcook04/plectis-erdos}
% BEGIN generated_semantic_coverage_macros
\newcommand{\SemanticAuthoredTheoremLike}{151,704}
\newcommand{\SemanticAuthoredInterpreted}{139,818}
\newcommand{\SemanticDirectEvidence}{3,330}
\newcommand{\SemanticContextual}{136,488}
\newcommand{\SemanticStructuralOnly}{11,886}
\newcommand{\SemanticAuthoredPercent}{92.2}
% END generated_semantic_coverage_macros
\newcommand{\repolink}[2]{\href{\repobase/blob/7ab6a901a2fe07e8d0e5c99dc46821adffc7411b/#1}{\texttt{#2}}}
\emergencystretch=1.7em

\runninghead{Problem-sized Lean worlds}
\title{Problem-Sized Lean Worlds}
\subtitle{Persistent, checkable research records for AI-assisted mathematics}
\author{Will Cook}
\date{September 2026}

\begin{document}
\maketitle
\paperprovenance

\begin{abstract}
Frontier models now produce candidate proofs of research problems faster than
mathematicians can read them. What is rarely preserved is the work around a
proof: the reasoning stays in a session, failed routes are discarded, and the
result reaches the public without the explanation, attribution and sense of
difficulty that make it usable. This paper describes a research environment
organised around a persistent unit of work, the problem-sized world: a set of
files in a public Git repository that keeps one open problem together with its
established results, failed routes, experiments, literature and exact open
obligations. A short paper compresses
its frontier for an expert; a long record keeps the complete arguments; Lean
checks the formal proofs; Comparator checks that each selected proof
establishes exactly its separately written statement; and a claim registry
records the wording, evidence class and limits of every public claim. Any model, provider or person can resume from the record,
and accepted work returns to it with credit. A public Lean repository maintains
eight open Erd\H{o}s problems this way. Its results include a kernel-checked
irrationality criterion for Problem 257 that admits supports with divergent
reciprocal sum, and a formal refutation of the Formal Conjectures statement of
Problem 1041, built on a counterexample posted by the forum contributor
\texttt{ani}; Formal Conjectures now records the answer as \texttt{False} and
links the proof. A comparison with Prove2Me, a hosted platform for stating and
proving Lean theorems, shows how the two designs compose.
\end{abstract}

\begin{paperresult}{Contribution and scope}
\textbf{Design.} The problem-sized world as the persistent unit of AI-assisted
research; a path from proof to public claim that separates what Lean settles,
what Comparator settles and what a mathematician judges; and a contribution
cycle that a person or an agent can run from a public clone.
\textbf{Instance.} Eight open Erd\H{o}s problems in one Lean repository:
sixteen problem papers totalling 773 pages, 1,823 Lean modules, and 674 results
asserted in the papers, of which 671 carry a Lean declaration and 633 have been
compared by Comparator.
\textbf{Scope.} For Problem 1041 the Formal Conjectures statement is refuted;
independent human review of its correspondence with the 1958 wording has not
been recorded. The other seven targets remain open. The system leaves novelty
and significance to experts.
\end{paperresult}

\section{What AI-assisted mathematics loses}\label{sec:intro}

Mathematicians writing in August and September 2026 describe one shift from
several sides. Kra writes that the machines are producing solutions
``faster than the mathematical community can read them, much less digest
them''~\cite{kra}. Cohn finds current models ``much better at solving
mathematical problems than at communicating the solutions well to humans'', and
calls the unexplained remainder technical debt: a correct result that someone
else must later pay to understand~\cite{cohn}. Tao argues that AI tools
flatten the difficulty landscape a field uses to choose its next questions,
that withheld negative results and processes hide where the new frontier lies,
and asks for analysis that ``identifies insights from the solution
process''~\cite{taomining}.
A declaration first signed by twenty-five Fields Medallists names what rushed
announcements lose: the writeup, the isolation of new methods and ideas, and
the citation of earlier work~\cite{fieldsdecl}. Sanderson compresses the point:
``every AI-generated proof is born an unsolved exposition
problem''~\cite{sanderson}.

Each of these describes a loss of state. A proof found in a session arrives
without the session: the routes that failed, the computations that chose the
route, the literature it relied on and the precise statement that remains open.
The next person, or the next model, begins again from the problem statement.
As producing results becomes easier, helping other people understand and use
them becomes the scarce part of the work.

This paper describes a research environment built around a persistent unit of
work, the \emph{problem-sized world}: one open problem together with everything
learned about it. A world is written for two readers at once. Its short paper
compresses the frontier so that an expert can inspect the underlying
mathematical ideas without reconstructing the whole search; its long record
keeps the complete arguments for whoever continues. Any model, provider or
person can resume from the world, so the cumulative mathematical work stays
available to be reasoned over and built on, and a stronger model starts from
the recorded frontier.

The paper makes four contributions: the problem-sized world as the persistent
unit of AI-assisted research (Section~\ref{sec:world}); a path from proof to
public claim that keeps three questions apart, namely whether Lean accepts a
proof, whether the proved statement is exactly the stated one, and whether that
statement means what the problem means (Section~\ref{sec:checks}); a
contribution cycle, published as agent skills in the public repository, through
which a person or an agent can take up an open obligation and return the result
with credit (Section~\ref{sec:cycle}); and a public instance on eight open
Erd\H{o}s problems (Section~\ref{sec:instance}), compared with Prove2Me and
related systems (Section~\ref{sec:related}).

\section{A problem-sized world}\label{sec:world}

A problem-sized world is a set of files in a public Git repository, organised
around the mathematics. Table~\ref{tab:world} lists its
parts.

\begin{table}[!ht]
\centering\small
\begin{tabular}{@{}>{\raggedright\arraybackslash}p{0.17\linewidth}
                    >{\raggedright\arraybackslash}p{0.50\linewidth}
                    >{\raggedright\arraybackslash}p{0.25\linewidth}@{}}
\toprule
\papertablehead
Part & What it holds & Written for \\
\midrule
Short paper & The problem and its status, the strongest results and the
mechanism behind them, relation to prior work, and the exact point where the
argument stops & An expert deciding whether to engage \\
\addlinespace
Long record & Complete arguments, failed routes, experiments with their finite
domains, corrected claims, and attribution for every borrowed idea & Whoever
continues the work \\
\addlinespace
Lean source & Formal statements and proofs checked by a pinned Lean~\cite{lean4} kernel &
Anyone who needs certainty about a formal statement \\
\addlinespace
Evidence marks & Beside each result, ``Lean'' links the declaration that
proves it, directly or through its evidence record, and ``Comparator'' links
its replay receipt; an unmarked result rests
on the paper's own argument & A reader checking one statement \\
\addlinespace
Claim registry & The registered wording of each public claim, with its
evidence class, the bounded domain of each finite instance and the open
statements it approaches & Release checks and anyone citing a claim \\
\addlinespace
Experiments & Code, inputs, finite domain, output and stated limits & Route
selection and replay \\
\bottomrule
\end{tabular}
\caption{The parts of one problem-sized world.}
\label{tab:world}
\end{table}

The short paper is the first read. It compresses the whole record, so an expert
can see the frontier of ideas in 12 to 28 pages and decide whether the problem
deserves their time. Its margin marks show, result by result, whether belief
rests on Lean, on Lean and Comparator, or on the argument printed in the paper;
a dagger on a Lean mark flags a named input the formal proof assumes. The long record holds what the essays above find missing
from announced results: the solution process, the routes that failed and why,
and the credit owed to earlier work. The difficulty landscape Tao
describes~\cite{taomining} is kept explicitly, as the set of results,
obstructions and open obligations around the endpoint.

\subsection{Failed routes and open obligations}\label{sec:mathloop}
\papersectiontarget{systems-mathloop}

The records treat negative results as outputs. A failed attempt shows that one
search path did not close. A counterexample refutes a conjecture on its exact
domain. A Lean no-go theorem rules out a class of strategies under explicit
hypotheses. The last two prevent repeated work and show which new ingredient an
endpoint needs. For Problem 257, a proposed estimate said that the greedy test
ruling out a rational number with denominator $Q$ as the sum of an infinite
subseries of $\sum_n 1/(2^n-1)$ stops within about $2\log_2 Q-3.3$ steps; an
\repolink{research/experiments/choices_contraction/README.md}{exact rational
probe} first rules out $189/388$ at step 17. The record keeps both that this
defeats the estimate and that it leaves the problem itself untouched.

\FloatBarrier
Open obligations are stated as exactly as results. Problem 249 asks whether
$S=\sum_{n\ge1}\varphi(n)/2^n$ is irrational, where $\varphi$ is Euler's
totient function~\cite{erdosgraham}. The development expresses irrationality
through finite certificates along $H_t=\operatorname{lcm}(1,\ldots,t)$. Writing
$\mathrm{Cert}(t)$ for the existence of a finite computation showing that a
scaled tail difference of $S$ at $H_t$ is not an integer, Lean checks
\[
\forall t\le82,\quad \mathrm{Cert}(t),
\]
and Lean also proves that irrationality is equivalent to the unbounded supply
\[
\forall T,\ \exists t>T,\quad \mathrm{Cert}(t).
\]
The finite theorem makes no \(t=83\) or cofinal claim. What remains open is the
unbounded supply itself, which is the irrationality of $S$ in equivalent form,
and the record states it in exactly that form.

\section{From a proof to a public claim}\label{sec:checks}
\papersectiontarget{systems-lifecycle}

A proof becomes a public claim by crossing three questions that the word
``verified'' tends to merge. Table~\ref{tab:questions} keeps them apart.

\begin{table}[!ht]
\centering\small
\begin{tabular}{@{}>{\raggedright\arraybackslash}p{0.36\linewidth}
                    >{\raggedright\arraybackslash}p{0.29\linewidth}
                    >{\raggedright\arraybackslash}p{0.27\linewidth}@{}}
\toprule
\papertablehead
Question & Answered by & Recorded in \\
\midrule
Does the proof establish this formal statement? & The pinned Lean kernel, in
continuous integration & Lean source and build receipts \\
\addlinespace
Is the proved statement exactly the separately stated one, using only
permitted axioms? & Comparator, with the Lean kernel and an independent checker
& The Comparator ledger and evidence records \\
\addlinespace
Does the formal statement say what the problem says, and does the result
matter? & A mathematician & Claim wording in the registry, and review events
when they occur \\
\bottomrule
\end{tabular}
\caption{Three questions a public mathematical claim must cross.}
\label{tab:questions}
\end{table}

Lean verifies that a proof establishes the formal statement written in the
source; it does not verify whether that statement captures the intended
mathematics or whether the paper describes it well. The gap is measurable.
Feng et al.\ had mathematicians grade 200 candidate solutions to Erd\H{o}s
problems produced by Aletheia, an agent built on Gemini Deep Think: 63 were
technically correct and 13 meaningfully correct, and the other 50 answered a
reading of the problem that missed Erd\H{o}s's intent~\cite{aletheia}.

Comparator, a checker the Lean FRO built for judging machine-written
proofs~\cite{leanfrocomparator}, answers the second question for selected statements.
A trusted challenge module restates each statement without its proof; the proof
must be a term of exactly that type, using only the permitted axioms, and the
Lean kernel must accept it. The corpus replay also passes every proof through
nanoda, an independent kernel~\cite{nanodalib}, and continuous integration requires
a deliberately altered statement to fail, so the harness cannot pass
vacuously. Comparator cannot decide whether a statement is the
right translation of the problem. It guarantees that the statement a reader
inspects is the statement that was proved, so ``Comparator-checked'' is
accurate where ``independently verified'' would overclaim.

The third question stays with people. The claim registry records the wording
of each public claim with its evidence class, the bounded domain of each finite
instance and the open statements it approaches, and the papers, README and
generated views may say only what the registry allows. In this repository the
maintainer is responsible for that wording, and no independent mathematical
review is yet recorded. A release checker keeps those surfaces consistent
with the registry and the Lean source, and requires a registered limitation to
stay visible wherever its claim appears.

That requirement comes from a test in which a limitation escaped. One
deliberately false README edit changed a clause saying that the finite cases of Problem 249 did not supply the
open requirement into one saying that they completed it. Lean was untouched,
and the checker passed because that relationship had not been registered.
After registration, the post-repair witness accepts the current README and
rejects a test copy containing the false clause. A small historical study
tested the checker with deliberately false edits: nine of the ten edits were
rejected. One escaped, for the reason just given. The edits were authored by the checker's author. The
original run logs were not retained. The other nine edits were not rerun
against the extended checklist.
The evidence marks a coverage boundary, not a reliability score.

The workflow does not technically force a second independent mathematician:
one maintainer can edit a Lean statement, its record and its prose together so
that every comparison agrees with the same mistake. Independent review is
recorded as an external event when it happens.

\section{The contribution cycle}\label{sec:cycle}

Someone outside the original research can contribute a correction or take an
argument further, and that work returns to the same record for that problem. The public
repository publishes this research, checking and revision cycle as thirteen
reusable agent skills that install into common agent harnesses. Each skill owns
one job: its starting state, the changes it may make, the evidence it must
return, its stopping condition and the next owner. Table~\ref{tab:cycle} shows
the ordinary path.

\papersectiontarget{systems-job-lifecycle}
\begin{table}[!ht]
\centering\small
\begin{tabular}{@{}>{\raggedright\arraybackslash}p{0.12\linewidth}
                    >{\raggedright\arraybackslash}p{0.33\linewidth}
                    >{\raggedright\arraybackslash}p{0.47\linewidth}@{}}
\toprule
\papertablehead
Step & Public skill & What it returns \\
\midrule
Orient & \texttt{explain-\allowbreak public-\allowbreak system},
\texttt{explore-\allowbreak the-\allowbreak corpus} & An account of the record
adapted to its reader, citing the evidence behind each statement \\
\addlinespace
Work & \texttt{mine-\allowbreak open-\allowbreak problem} & The smallest
evidence-bearing change: a theorem, a counterexample, a no-go result or a
diagnosed failure \\
\addlinespace
Validate & \texttt{lean-\allowbreak concurrent-\allowbreak validation},
\texttt{land-\allowbreak lean-\allowbreak proofs} & One Lean verdict per
distinct request, and every paper that states the result linked to its
declaration \\
\addlinespace
Propagate & \texttt{propagate-\allowbreak research-\allowbreak consequences} &
A recorded disposition for every affected import, claim, paper and open
obligation \\
\addlinespace
Package & \texttt{erdos-\allowbreak research-\allowbreak return},
\texttt{public-\allowbreak mathematical-\allowbreak writing} & A
provenance-preserving return, with its prose held to the evidence \\
\addlinespace
Propose & \texttt{submit-\allowbreak pull-\allowbreak request} & A pull request
from a named starting commit \\
\addlinespace
Adopt & Maintainer review under the credit policy & Reconciliation with the
current tree and a public credit receipt \\
\bottomrule
\end{tabular}
\caption{The contribution cycle as published in the public repository. Four
further skills coordinate discovery and stewardship, admit a new problem world,
maintain the clone's own tooling and install the set. A cycle closes when the
change has evidence and every consequence has a recorded disposition; the open
problem closes only with a proof of its original statement.}
\label{tab:cycle}
\end{table}

Two steps carry most of the discipline. Propagation starts from the changed
object and lists every reverse import, claim record, paper, experiment and open
obligation the change may affect. Each must be updated, verified unchanged,
deferred with a reason or marked out of scope, and an empty search never counts
as evidence of no consequence. Adoption replays the contributor's original
change from its named starting commit before reconciling it with the current
tree, so the original work and any conflict resolution are credited separately.
Orientation needs no knowledge of the file layout: \texttt{explain-public-system}
adapts its account to a lay reader, a mathematician, a formaliser, a compute
contributor or an infrastructure contributor. The
\papersectionlink{open-source-mathematics-strategy.pdf}%
{strategy-protocol}{companion contribution protocol} gives the full sequence.

\paragraph{Discovery and stewardship.}\papersectiontarget{systems-coupled-goals}
Long-running work uses two coupled roles. A discovery pass stays close to the
frontier: it runs discriminating computations, attempts proofs and returns
results. A stewardship pass reads each new result against the whole corpus and
decides what it establishes, whether it is routine or strong, where it belongs
in a paper, and which obligation deserves the next unit of work. One agent may
play both roles in sequence. Neither inherits the other's authority, and both
wait while the frontier and its consumers are unchanged.

\paragraph{Credit.}
The \repolink{docs/research-commons/CREDIT_POLICY.md}{credit policy} lets an
idea be credited separately from its proof, its Lean formalisation and its
explanation, the roles Kra asks journals to distinguish~\cite{kra}, and credits
software as a role of its own. A first
return can be a plain-language research-progress issue or
email: no clone, proof, code or receipt schema is required. Accepted work
receives a public receipt tied to exact artifacts. Acceptance does not
establish theorem status, novelty, or release inclusion. Corrections append to
the history, and an earlier contributor's record stays in it.

\subsection{Worked example: Problem 1041}\label{sec:example}

Problem 1041 shows the cycle end to end, with the decisive idea arriving from
outside. It asks whether, for a monic polynomial $f$ with all roots in the open
unit disc, two roots can always be joined inside the lemniscate $\{|f|<1\}$ by
a path of length less than~$2$~\cite{erdos1041,ehp1958}. On 7 September 2026
the erdosproblems.com contributor \texttt{ani} posted a degree-seven
counterexample, found, as the post says, with the help of GPT-6~\cite{aniforum},
and the record took the construction as an attributed input five days later.
Lean first checked the total-variation form. On 23 September it proved that for
that polynomial every preconnected subset of the strict lemniscate containing
two distinct roots has one-dimensional Hausdorff measure greater than~$2$,
which refutes the exact path-image-length statement in Google DeepMind's Formal
Conjectures repository~\cite{formalconjectures}. The same day, the Formal
Conjectures maintainers merged a change that marks the problem solved with
answer \texttt{False} and links this proof~\cite{fcpr}. The short paper and long
record carry the proof with its Lean mark and credit \texttt{ani} for the
construction. One judgement remains open: whether the formal statement matches
the 1958 wording~\cite[Problem~5, p.~139]{ehp1958}. The record says so beside
the result.

\section{Eight problems in one repository}\label{sec:instance}
\papersectiontarget{systems-public}

The \href{\repobase}{public repository} maintains eight open Erd\H{o}s
problems~\cite{erdosproblems} as problem-sized worlds. Its README leads with
Problem 257, which asks whether $\sum_{a\in A}(2^a-1)^{-1}$ is irrational for
every infinite set $A$ of positive integers. For a set $A$ of positive integers
and a finite nonempty set $P$ of primes, let $h(a)=\prod_{p\in P}p^{v_p(a)}$ be
the $P$-part of $a$, where $v_p(a)$ is the exponent of $p$ in $a$. If
\[
\sum_{a\in A}\frac{h(a)}{a\,(2^{h(a)}-1)}<\infty,
\]
then $\sum_{a\in B}(b^a-1)^{-1}$ is irrational for every integer $b\ge2$ and
every infinite $B\subseteq A$. The condition admits sets $A$ whose reciprocal
sum diverges, and Lean checks the criterion. Table~\ref{tab:problems} gives one
result for each problem and the point where it stops.

\begin{table}[!ht]
\centering\small
\begin{tabular}{@{}>{\raggedright\arraybackslash}p{0.17\linewidth}
                    >{\raggedright\arraybackslash}p{0.50\linewidth}
                    >{\raggedright\arraybackslash}p{0.25\linewidth}@{}}
\toprule
\papertablehead
Problem & A result to start with & Where it stops \\
\midrule
\#68, factorial denominator & Lean checks that the logarithm of the uncleared
common denominator grows at least like $N^{3/2}\log N$, and that irrationality
is equivalent to cofinally many non-unit factorial carries & Those carries are
not produced \\
\addlinespace
\#243, Sylvester-tail rigidity & Lean checks irrationality of $\sum 1/a_n$ under the
exact cubic rate $a_n^2/a_{n+1}=1+3/n+o(n^{-3})$ & The unrestricted
Sylvester-tail question \\
\addlinespace
\#249, binary totient series & In every base $k$, the totient sections
$\varphi(k^jn+r)$ with $j\le e$ span dimension $k^e+1$; certificates for every
$t\le82$ & Certificates for arbitrarily large $t$ \\
\addlinespace
\#251, prime-gap dyadic series & Lean checks a synthetic prime-gap countermodel; the paper
proves a sparse-perturbation obstruction & An argument specific to actual
prime gaps \\
\addlinespace
\#257, infinite exponent supports & The weighted-support criterion above, at
every integer base & Arbitrary infinite supports \\
\addlinespace
\#269, three-prime running LCMs & Both two-prime sums are transcendental, by a cited
Hecke--Mahler theorem, the repeated-sum case first posted by Fan~\cite{fan269};
Lean checks the formulas and the conditional transfer & The three-prime case \\
\addlinespace
\#1041, lemniscate connections & The Formal Conjectures statement is refuted with
\texttt{ani}'s example & Match with the 1958 wording \\
\addlinespace
\#1049, rational-base Lambert series & Lean checks irrationality in Zudilin's rational-base
region and exact Hankel orders & Base $3/2$ and all rational bases \\
\bottomrule
\end{tabular}
\caption{One result per problem, with its stopping point. The papers state
each result in full with its hypotheses and sources.}
\label{tab:problems}
\end{table}

The repository holds 1,823 Lean modules under a pinned
toolchain, sixteen problem papers totalling 773 pages, a cross-problem
synthesis paper and three systems papers. The papers assert 674 results. Of
these, 671 carry a Lean mark, linking a declaration that proves the result
exactly, in a stronger form or from a named input, and 633 have been compared
by Comparator; every receipt in the pinned corpus replay was accepted by both
the Lean kernel and nanoda. The claim registry holds 150 public claims in seven
statuses, from proved here and formalised here to conditional reduction and
open, with 19 open propositions stated exactly. Continuous integration builds
the supported Lean roots, a coverage build compiles every module a paper cites,
and the release checker runs 15,666 checks over claims, papers, generated views
and licences. A fresh clone needs no private file, and no public theorem
depends on an unpublished lemma.

\section{How the records are produced}\label{sec:production}
\papersectiontarget{systems-comprehension}

One person has built, over the past year, the private environment that produces
the records, and runs agents from several model providers in it, including
Anthropic's Claude and OpenAI's GPT models, concurrently on one shared file
system. Five design choices carry most of the
weight.

\paragraph{Shared state on disk.} Durable state lives in files: sources,
structured records, append-only ledgers, validation receipts and authored
explanations. A change exists once its file is written and its
checker accepts it, so work resumes across sessions, models and providers.

\paragraph{Bounded context.} Before an agent opens a large source, a router
compiles a packet conditioned on its task: the endpoint, its status and claim
ceiling, the strongest nearby results, known falsifiers and literature, each
labelled by evidence class, with a list of what was omitted and the command
that restores it. A packet that would exceed its budget reports the overflow
and stops.

\paragraph{Claims before writes.} An agent claims exact paths for a bounded
lease before changing them. Overlapping claims wait or coordinate, expired
leases lose authority, and parallel work meets at a fan-in step where results
are compared and validated.

\paragraph{Single-flight validation.} Lean builds pass through a semantic
single-flight queue keyed by source, targets, toolchain and dependency lock.
Equivalent requests share one build. An agent that cannot get the build slot
receives a deferral, recorded separately from theorem failures, and continues
with independent work.

\paragraph{Route repair.} A failure in navigation, scheduling,
experimentation or validation can change a route, skill or check for later
agents through a guarded proposal. It never changes the status of a
mathematical claim.

The public repository replays without the private environment. The private
environment explains how the records were produced; the public evidence is
what supports each claim.

\section{Relation to Prove2Me and other systems}\label{sec:related}

Prove2Me~\cite{prove2me} is the closest published system. It is a hosted
platform on which anyone with a coding agent states Lean theorems and submits
proofs. Each statement is immutable and may carry several proofs; the server
accepts a proof only when its type is exactly the target's, with no
\texttt{sorry} and only whitelisted axioms, and a disproof proves the negation.
Missions on papers, textbooks and open problems confine human audit to a goal,
its definitions and milestone lemmas, aided by a blind agent read-back of each
Lean statement, and a proof sketch may import open statements for other agents
to close. Prove2Me began as a class game at Columbia in spring 2026 and its
paper appeared on 28 August 2026; the first public version of this paper, with
its claim registry and release checks, appeared in July 2026, and the two
designs developed independently. Prove2Me is the published reference for hosted
statement and proof separation, audited mission cores and read-back auditing.
Table~\ref{tab:prove2me} sets the two side by side.

\begin{table}[!ht]
\centering\small
\begin{tabular}{@{}>{\raggedright\arraybackslash}p{0.19\linewidth}
                    >{\raggedright\arraybackslash}p{0.36\linewidth}
                    >{\raggedright\arraybackslash}p{0.37\linewidth}@{}}
\toprule
\papertablehead
 & Prove2Me & Problem-sized worlds \\
\midrule
Unit of work & One immutable statement or proof within a mission & One open
problem with its whole record \\
\addlinespace
What persists & Statements, proofs, sketches, explanations, milestones and
discussion; disproofs and deprecations stay visible & Short and long papers,
failed routes, experiments, typed claims and open obligations beside the Lean
source \\
\addlinespace
Exact-statement check & Every submission, on the server & Selected statements,
through Comparator, with a second kernel in the corpus replay \\
\addlinespace
Statement fidelity & A captain audits each mission core with a blind read-back;
a moderator approves public missions & The maintainer is responsible for the
wording of each public claim in the claim registry; no independent review is
recorded yet \\
\addlinespace
Exposition & A description for each statement and an explanation for each proof
& A short paper and a long record for each problem, with evidence marks at
each result \\
\addlinespace
Credit & Permanent names, trust scores, leaderboards and import citations &
Receipts by contribution role, and attribution for outside ideas \\
\addlinespace
Contribution path & Web platform and API, with an account & Git clone and
thirteen agent skills, a pull request, an issue or an email \\
\addlinespace
Scale, 24 September 2026 & 373 missions, 88,246 theorems, 722 users & Eight
problems, 1,823 Lean modules, one maintainer \\
\bottomrule
\end{tabular}
\caption{Prove2Me and problem-sized worlds, from Prove2Me's paper and public
pages and from this repository, on 24 September 2026.}
\label{tab:prove2me}
\end{table}

The two compose. Prove2Me contributes hosting, importable statements,
per-submission checking and a contributor base; a world contributes the
problem-level context around a statement: its papers, failed routes,
obstructions and open obligations. The Problem 1041 package from this repository
is hosted on Prove2Me, where its 21 theorems have accepted proofs and
\texttt{ani} is credited separately for the construction.

\paragraph{Other systems.} Feng et al.\ ran an agent over the 700 Erd\H{o}s
problems then listed as open and published graded outcomes~\cite{aletheia}, and
propose reporting AI-assisted results on at least two axes, autonomy and
significance, with significance judged only by mathematicians expert in the
area~\cite{autonomousmath}; the claim registry records the evidence behind a
result and leaves its significance to experts in the same way. Tao's community wiki listed AI contributions to the Erd\H{o}s problems,
failed attempts included, until its updates stopped on 30 June
2026~\cite{taowiki}; a registered claim, with its evidence class and limits, is
the kind of record such a list can cite. Formal
Conjectures~\cite{formalconjectures} turns open problems into formal statements
that automated systems can target. Ringer distinguishes mathematics as a
process from mathematics as a benchmark~\cite{ringer}; a world keeps the
process around such a statement. Blueprints connect informal proof plans to
Lean declarations~\cite{leanblueprint,leanarchitect}, and provers such as
ReProver, built on the LeanDojo toolkit, and OpenProver~\cite{leandojo,openprover}
search for proofs; a world can use a blueprint as its dependency map and any
such prover as a source of candidate proofs.
Li et al.\ published a human-in-the-loop protocol for research mathematics that
holds generation and validation apart without a proof assistant~\cite{lihai}.

\section{Limits, and what stronger models change}\label{sec:limits}
\papersectiontarget{systems-trust}

The other seven targets remain open, and independent human review of the
correspondence between the Problem 1041 refutation and the 1958 wording has not
been recorded. The system records evidence classes and review events;
novelty and significance are judgements it leaves to experts, as Feng et al.\
also propose~\cite{autonomousmath}. One maintainer can make a coordinated
mistake that every structural check accepts. By 24 September 2026 no outside
contributor had opened a pull request or issue in the repository; its outside
record is in Formal Conjectures, whose maintainers have merged four of its
contributions, and on Prove2Me, which hosts its Problem 1041 package with
accepted proofs. A record also leaves several of the concerns in
Section~\ref{sec:intro} where it finds them: it trains no one, it restores no
signal of individual expertise~\cite{litt}, and it answers neither the ethical
objections some mathematicians raise to AI-assisted mathematics~\cite{chu} nor
the warning of forty-two Fellows and Foreign Members of the Royal Society about
catastrophic risk from AI~\cite{greenletter}. The private production
environment is described here and is not distributed.

\papersectiontarget{systems-scaling}
The same design serves stronger models. A stronger model starts from the
current frontier with the failed routes already marked,
and it can reorganise, re-explain and extend the record as well as the
mathematics.
The public interfaces admit any producer: an outside contributor can attach any
agent runner to the clone, and a laboratory with many frontier models could
replace the private production environment while keeping the same evidence
boundary.

Gowers suggests a well-designed database of what is known, probably built with
AI help, and Koukoulopoulos a public facility in which coordinated agents
explore a research programme divided into subproblems under researchers'
direction~\cite{gowers,koukoulopoulos}. A problem-sized world supplies the
shared record both proposals need, at the scale of eight problems. Antieau
writes that texts generated largely by a language model deserve a home other
than the arXiv~\cite{antieau}; a world gives its long records one beside the
proofs they explain. The same structure suits a single landmark result: a short
paper explaining its core ideas and relation to prior work, connected to the
full argument, the relevant formal checks and the questions it opens, with
outside corrections and extensions returning to the same record. The design
asks three things of a field: stable result identities, a checker where one
exists, and an explicit boundary on public claims. The next measurement is
comparative: whether agents given a world choose better frontiers, repeat fewer
dead ends and reach useful obstructions sooner than agents given the same
mathematics as prose.

\section{Conclusion}\label{sec:conclusion}

A problem-sized world keeps the record around a proof that lets other people
understand it, check it, credit it and build on it: its frontier, failed routes,
formal checks, attribution and exact open statements. The eight worlds in the
public repository show that one person can keep that record persistent,
checkable and open to contribution across several hard problems. The record is
built for models stronger than the ones that produced it: each starts from the
recorded frontier with the failed routes marked, and
Section~\ref{sec:limits} names the measurement that tests whether this makes
it more effective.

\appendix
\section{Inspection routes and reproducibility}\label{app:repro}

\begin{table}[H]
\centering\small
\begin{tabular}{@{}>{\raggedright\arraybackslash}p{0.36\linewidth}
                    >{\raggedright\arraybackslash}p{0.56\linewidth}@{}}
\toprule
\papertablehead
Question & Start here \\
\midrule
How does the repository fit together? &
\repolink{docs/ARCHITECTURE.md}{docs/ARCHITECTURE.md} \\
What is proved and what is open? &
\repolink{docs/RESULTS.md}{docs/RESULTS.md} and
\repolink{docs/claims.json}{docs/claims.json} \\
Where is the checked mathematics? &
\repolink{lean/Erdos249257.lean}{lean/Erdos249257.lean} and
\repolink{lean/ErdosProblems.lean}{lean/ErdosProblems.lean} \\
What does Comparator check? &
\repolink{docs/EXTERNAL_VERIFICATION.md}{docs/EXTERNAL\_VERIFICATION.md} and
\repolink{verification/comparator.json}{verification/comparator.json} \\
Which checks gate a release? &
\repolink{scripts/check_release.py}{scripts/check\_release.py} and
\repolink{.github/workflows/lean.yml}{.github/workflows/lean.yml} \\
How can work return with credit? &
\repolink{CONTRIBUTING.md}{CONTRIBUTING.md},
\repolink{docs/research-commons/CREDIT_POLICY.md}{credit policy} and
\repolink{docs/research-commons/ARCHITECTURE_CONTRIBUTIONS.md}{architecture contributions} \\
Which agent skills exist? &
\repolink{skills/README.md}{skills/README.md} \\
\bottomrule
\end{tabular}
\caption{Entry points for inspection. They lead to the authority-bearing files
and add no authority of their own.}
\label{tab:routes}
\end{table}

A fresh clone reproduces the control card and the structural checks without
Lean:
\begin{verbatim}
python3 scripts/proof_cockpit.py --format card
python3 scripts/proof_cockpit.py --check
python3 scripts/check_release.py
\end{verbatim}
Formal authority begins with the pinned \texttt{lake build} named by the card.
The paper inventory, \path{docs/publication_contract.json}, records source and
PDF hashes and validation commands, and the evidence for the checker study in
Section~\ref{sec:checks} is recorded in \path{docs/publication_evidence.json}.

{\footnotesize
\begin{multicols}{2}
\begin{thebibliography}{99}
\setlength{\itemsep}{1pt}
\setlength{\parsep}{0pt}
\bibitem{kra} B.~Kra, \emph{Deep theorems were scarce and difficult and so
became an effective mechanism to identify deep thought. AI has broken this
system}, guest post on \emph{What's new}, 13 September 2026,
\href{https://terrytao.wordpress.com/2026/09/13/deep-theorems-were-scarce-and-difficult-and-so-became-an-effective-mechanism-to-identify-deep-thought-ai-has-broken-this-system/}{blog post}.
\bibitem{cohn} H.~Cohn, \emph{The technical debt of AI-generated mathematics},
guest post on \emph{What's new}, 15 September 2026,
\href{https://terrytao.wordpress.com/2026/09/15/the-technical-debt-of-ai-generated-mathematics/}{blog post}.
\bibitem{taomining} T.~Tao, thread on mining open problems, Mathstodon,
8 September 2026, \href{https://mathstodon.xyz/@tao/117237320796901560}{thread}.
\bibitem{fieldsdecl} Twenty-five Fields Medallists, \emph{A severe
misalignment of AI in mathematics}, declaration posted on \emph{What's new},
11 September 2026,
\href{https://terrytao.wordpress.com/2026/09/11/a-severe-misalignment-of-ai-in-mathematics/}{blog post}.
\bibitem{sanderson} G.~Sanderson, \emph{If math is more than proof, we need to
better celebrate the rest of it}, guest post on \emph{What's new}, 18 September
2026, \href{https://terrytao.wordpress.com/2026/09/18/if-math-is-more-than-proof-we-need-to-better-celebrate-the-rest-of-it/}{blog post}.
\bibitem{erdosgraham} P.~Erd\H{o}s and R.~L.~Graham, \emph{Old and New
Problems and Results in Combinatorial Number Theory}, Monographies de
L'Enseignement Math\'ematique 28, 1980, p.~61.
\bibitem{lean4} L.~de Moura and S.~Ullrich, \emph{The Lean 4 Theorem Prover
and Programming Language}, in \emph{Automated Deduction, CADE 28}, Lecture
Notes in Computer Science 12699, 2021, pp.~625--635,
\href{https://doi.org/10.1007/978-3-030-79876-5_37}{DOI}.
\bibitem{aletheia} T.~Feng, T.~Trinh, G.~Bingham, et al., \emph{Semi-Autonomous
Mathematics Discovery with Gemini: A Case Study on the Erd\H{o}s Problems},
2026, \href{https://doi.org/10.48550/arXiv.2601.22401}{arXiv}.
\bibitem{leanfrocomparator} Lean FRO, \emph{Comparator}, 2025,
\href{https://github.com/leanprover/comparator}{GitHub}.
\bibitem{nanodalib} \texttt{ammkrn}, \emph{nanoda\_lib}, an independent type
checker for Lean 4, \href{https://github.com/ammkrn/nanoda_lib}{GitHub}.
\bibitem{erdos1041} T.~F.~Bloom, \emph{Erd\H{o}s Problem \#1041},
\href{https://www.erdosproblems.com/1041}{erdosproblems.com}, accessed
September 2026.
\bibitem{ehp1958} P.~Erd\H{o}s, F.~Herzog and G.~Piranian, \emph{Metric
properties of polynomials}, J. Analyse Math. 6 (1958), 125--148,
\href{https://doi.org/10.1007/BF02790232}{DOI}.
\bibitem{aniforum} \texttt{ani}, post in the Problem 1041 discussion thread,
7 September 2026,
\href{https://www.erdosproblems.com/forum/thread/1041\#post-8861}{erdosproblems.com forum}.
\bibitem{formalconjectures} Google DeepMind, \emph{Formal Conjectures},
\href{https://github.com/google-deepmind/formal-conjectures}{GitHub}, accessed
September 2026.
\bibitem{fcpr} Formal Conjectures, pull request 6505, \emph{Erd\H{o}s 1041:
mark solved with answer(False) and link a formal proof}, merged 23 September
2026, \href{https://github.com/google-deepmind/formal-conjectures/pull/6505}{GitHub}.
\bibitem{erdosproblems} T.~F.~Bloom, \emph{Erd\H{o}s problems},
\href{https://www.erdosproblems.com}{erdosproblems.com}, accessed September 2026.
\bibitem{fan269} S.~Fan, comment on Erd\H{o}s Problem \#269, erdosproblems.com
forum, 26 June 2026,
\href{https://www.erdosproblems.com/forum/thread/269\#post-7218}{forum post}.
\bibitem{prove2me} S.~Chen, K.~Marwaha, X.~Lu, H.~Yuen, and T.~Peng,
\emph{Prove2Me: An Open Collaborative Platform for Scaling Math
Formalization}, 2026, \href{https://doi.org/10.48550/arXiv.2608.28433}{arXiv}.
\bibitem{autonomousmath} T.~Feng, T.~H.~Trinh, G.~Bingham, et al.,
\emph{Towards Autonomous Mathematics Research}, 2026,
\href{https://doi.org/10.48550/arXiv.2602.10177}{arXiv}.
\bibitem{taowiki} T.~Tao, \emph{AI contributions to Erd\H{o}s problems},
\href{https://github.com/teorth/erdosproblems/wiki/AI-contributions-to-Erd%C5%91s-problems}{GitHub},
accessed September 2026.
\bibitem{ringer} T.~Ringer, \emph{Becoming a benchmark}, guest post on
\emph{What's new}, 17 September 2026,
\href{https://terrytao.wordpress.com/2026/09/17/becoming-a-benchmark/}{blog post}.
\bibitem{leanblueprint} P.~Massot, \emph{leanblueprint}, plasTeX plugin for
Lean formalisation blueprints, 2020,
\href{https://github.com/PatrickMassot/leanblueprint}{software repository}.
\bibitem{leanarchitect} T.~Zhu, P.~Monticone, S.~Welleck, and J.~Avigad,
\emph{LeanArchitect: Automating Blueprint Generation for Humans and AI},
in \emph{17th International Conference on Interactive Theorem Proving},
LIPIcs 382, 2026, pp.~25:1--25:16.
\bibitem{leandojo} K.~Yang et al., \emph{LeanDojo: Theorem Proving with
Retrieval-Augmented Language Models}, NeurIPS 2023.
\bibitem{openprover} M.~Kripner and M.~Straka, \emph{OpenProver: Agentic and
Interactive Theorem Proving with Lean 4}, 2026,
\href{https://doi.org/10.48550/arXiv.2607.09217}{arXiv}.
\bibitem{lihai} C.~Li, Z.~Lai, D.~An, J.~Hu, and Z.~Wen, \emph{Advancing
Mathematical Research via Human-AI Interactive Theorem Proving}, 2025,
\href{https://arxiv.org/abs/2512.09443v2}{arXiv:2512.09443v2}.
\bibitem{litt} D.~Litt, \emph{A beginning for mathematics}, 13 September
2026, \href{https://www.daniellitt.com/blog/2026/9/13/a-beginning-for-mathematics/}{blog post}.
\bibitem{chu} T.~Chu, \emph{The AI dissenter viewpoint}, \emph{Proofs and
Prompts}, 9 August 2026,
\href{https://proofsandprompts.com/2026/08/09/the-ai-dissenter-viewpoint/}{blog post}.
\bibitem{greenletter} B.~Green and forty-one other Fellows and Foreign Members
of the Royal Society,
\emph{Open letter to Sir Paul Nurse, President of the Royal Society},
\emph{Proofs and Prompts}, 17 September 2026,
\href{https://proofsandprompts.com/2026/09/17/open-letter-to-sir-paul-nurse-president-of-the-royal-society/}{blog post}.
\bibitem{gowers} W.~T.~Gowers, \emph{Why I didn't sign the Fields medallists'
letter}, 17 September 2026,
\href{https://gowers.wordpress.com/2026/09/17/why-i-didnt-sign-the-fields-medallists-letter/}{blog post}.
\bibitem{koukoulopoulos} D.~Koukoulopoulos, \emph{A CERN for AI-assisted
science?}, guest post on \emph{What's new}, 17 September 2026,
\href{https://terrytao.wordpress.com/2026/09/17/a-cern-for-ai-assisted-science/}{blog post}.
\bibitem{antieau} B.~Antieau, \emph{Fast math/slow math}, 15 September 2026,
\href{https://antieau.github.io/2026/09/15/fast-math-slow-math.html}{blog post}.
\end{thebibliography}
\end{multicols}
}
\end{document}
